% !TEX program = xelatex
\documentclass[11pt,a4paper]{ctexart}
\usepackage[margin=1.5cm]{geometry}
\usepackage{enumitem}
\usepackage[hidelinks]{hyperref}
\setlength{\parindent}{0pt}
\setlength{\parskip}{4pt}
\pagenumbering{gobble}
\newcommand{\head}[1]{\vspace{2pt}\textbf{\large #1}\vspace{2pt}\hrule\vspace{4pt}}

\begin{document}

\begin{center}
{\LARGE \textbf{方保骅\;\;(Arno Fang)}}\\
深圳 / 上海 \quad
\href{mailto:baohuafang@link.cuhk.edu.cn}{baohuafang@link.cuhk.edu.cn} \quad
+86 150 2689 3718 \\
\href{https://hiddenblade.net}{hiddenblade.net} \quad
\href{https://github.com/Hidden61ade}{github.com/Hidden61ade} \quad
\href{https://hidden61ade.itch.io/}{hidden61ade.itch.io}
\end{center}

\head{求职目标}
散爆网络｜关卡策划实习生

\head{教育背景}
\textbf{香港中文大学（深圳）} \hfill 2023.09 -- 至今\\
电子与计算机工程（Computer Engineering）本科 \hfill GPA: \textbf{3.50/4.00（预计）}

\head{岗位匹配亮点}
\begin{itemize}[leftmargin=1.3em,itemsep=2pt,topsep=2pt]
  \item \textbf{深度 FPS/TPS 玩家：}FPS/TPS 累计 1000+ 小时，长期体验 Apex、CS、COD、R6、泰坦陨落2、The Finals、瓦洛兰特等主流射击产品。
  \item \textbf{关卡与机制实现能力：}具备 Unity 开发经验，能完成玩法原型、技能逻辑、关卡触发与调试；UE5 可进行基础操作。
  \item \textbf{3C 体验理解：}持续进行射击品类竞品拆解，关注角色移动反馈、镜头手感、技能可读性与战斗节奏。
  \item \textbf{团队协作与推进：}多次 Game Jam 担任 Producer/带队，含 10 人项目，具备任务拆解、版本推进与跨职能沟通能力。
\end{itemize}

\head{实习与项目经历}
\textbf{腾讯科技（深圳）Tiki Studio｜游戏设计实习生} \hfill 2025\\
\begin{itemize}[leftmargin=1.3em,itemsep=2pt,topsep=2pt]
  \item 参与项目需求实现与迭代流程，协助设计内容快速原型验证并推动跨团队协作。
  \item 使用 Python 与 AI 工具优化资源生产流程，支撑高频版本迭代。
  \item 在“策划需求-技术实现-版本验证”闭环中沉淀协作经验，可直接迁移到关卡功能跟进与问题回归流程。
\end{itemize}

\textbf{Devil, Cops, Androids（Lead Designer / Programmer, Unity）} \hfill 2026\\
\textit{链接：\href{https://synapse75.itch.io/devilcopsandroids}{Itch.io}}\\
\begin{itemize}[leftmargin=1.3em,itemsep=2pt,topsep=2pt]
  \item Global Game Jam 2026 射击向项目，在 48 小时内完成核心战斗循环、关卡逻辑与 UI 流程搭建。
  \item 负责顶视角射击玩法原型实现与参数调试，围绕移动速度、射击反馈、敌我信息可读性进行快速迭代。
  \item 结合测试反馈定位卡关与节奏断点，迭代出生存压力与资源补给更平衡的关卡节奏。
\end{itemize}

\textbf{Heart Keys（导演/程序，10 人团队）} \hfill 2025\\
\textit{链接：\href{https://ldjam.com/events/ludum-dare/58/heart-keys}{LDJam}}\\
\begin{itemize}[leftmargin=1.3em,itemsep=2pt,topsep=2pt]
  \item 负责团队协作流程、版本控制与玩法联调，项目在 Ludum Dare 58 排名 \textbf{19/1083}（32 条评分）。
  \item 关卡采用“能力回收”成长线：基础移动到 Dash/Recall/Create 能力逐步解锁，驱动谜题层级递进。
  \item 根据玩家反馈持续修复关键流程问题，发布 1.1.2 版本解决主线钥匙无法获取与关卡卡死问题。
  \item 推进“设计-程序-美术”闭环，积累多人项目的关卡联调、可读性优化与测试回归经验。
\end{itemize}

\textbf{The\_Birthday\_Party（主策划/程序/编剧）} \hfill 2024\\
\textit{链接：\href{https://hidden61ade.itch.io/the-birthday-party}{Itch.io}}\\
\begin{itemize}[leftmargin=1.3em,itemsep=2pt,topsep=2pt]
  \item 完成 16000 字分支剧本与交互逻辑实现，将“信息获取-选择反馈-状态变化”做成可执行玩法链路。
  \item 通过脚本与系统联动实现元叙事交互，强化了“机制表达情绪”的关卡与节奏设计能力。
\end{itemize}

\textbf{Right Click To Activate Translator（内容与系统策划）} \hfill 2024 -- 至今\\
\textit{链接：\href{https://store.steampowered.com/app/3128590/_/}{Steam}}\\
\begin{itemize}[leftmargin=1.3em,itemsep=2pt,topsep=2pt]
  \item 像素横版语言解谜游戏，围绕“翻译器词典编辑-场景线索读取-谜题求解”构建核心循环。
  \item 负责语言系统与关卡信息投放设计，控制提示密度与探索自由度，平衡“可解性”与“发现感”。
  \item 项目获 2024 CUSGA 最佳解谜游戏，已上线 Steam Demo 页面用于持续收集外部反馈。
\end{itemize}

\textbf{GugugagaPenguin!（核心程序/交互原型）} \hfill 2026\\
\textit{链接：\href{https://www.bilibili.com/video/BV1DL9CBzEej/}{Bilibili 演示视频}}\\
\begin{itemize}[leftmargin=1.3em,itemsep=2pt,topsep=2pt]
  \item 基于 Leap Motion 2 的手势交互原型，采用“扇手飞行”作为核心输入映射，完成输入识别与反馈调优。
  \item 设计并实现收集硬币、打飞镖等轻量关卡目标，验证非传统输入设备在关卡交互中的可行性。
\end{itemize}

\textbf{Hashmon（全栈开发）} \hfill 2026\\
\textit{链接：\href{https://github.com/Hidden61ade/Hashmon}{GitHub}}\\
\begin{itemize}[leftmargin=1.3em,itemsep=2pt,topsep=2pt]
  \item 完成 Web3 游戏原型闭环：钱包连接、IPFS 元数据、Sepolia Mint、Inventory、Marketplace List/Buy。
  \item 通过 Battle/Garden 双场景复用同一 Companion 数据层，强化了“系统配置一致性 + 场景差异化体验”设计能力。
  \item 在课程级项目中承担需求拆解、功能落地与文档交接，提升了复杂系统下的排错与协同效率。
\end{itemize}

\head{竞品分析与玩家经历（射击向重点）}
\textbf{近期重点体验：}Marathon: 失落星船、逃离鸭可夫、Death Stranding 2、Deadlock。

\textbf{FPS/TPS（1000h+）：}
Apex Legends、战地系列、CS 系列、COD 系列、彩虹六号：围攻、泰坦陨落2、The Finals、瓦洛兰特。

\textbf{电竞赛事关注：}
自 2024 年起持续追踪 VCT 与 CS2 赛事，关注版本改动对阵容选择、战术节奏与观赛叙事的影响。

\textbf{其他类型（用于横向比较）：}
RPG（艾尔登法环、博德之门3、赛博朋克2077 等）、竞速（极限竞速、极品飞车、神力科莎）、主机（NS/3DS、Xbox One）、手游（王者荣耀、明日方舟等）。

\head{技能}
\textbf{引擎与工具：}Unity（熟练）、UE5（基础）、Blender（角色与动画基础）、Excel、Visio、Figma\\
\textbf{编程：}C\#、Python、C++、JavaScript\\
\textbf{能力：}关卡逻辑实现与调试、玩法原型、竞品拆解、跨团队沟通\\
\textbf{语言：}中文（母语），英文（TOEFL 110 / GRE 328）

\end{document}
